\begin{problem}[Squarefree hypersurfaces]
Let $R=k[x_1,\dots,x_n]$, a UFD, and let $0\ne f\in R$.  Prove that
\[
R/(f)\text{ is reduced}
\quad\Longleftrightarrow\quad
f\text{ is squarefree up to a unit}.
\]
Apply the result to $xy$, $x^2-y^3$, and $x^2y$.
\end{problem}
\paragraph{Why this problem matters.}
For principal equations, multiplicity of irreducible factors is exactly nilpotent thickness.
\paragraph{Useful ideas before starting.}
In a UFD, $(f)$ is radical exactly when no irreducible factor appears with exponent $>1$.
\paragraph{Guided hints.}
If $p^2\mid f$, find an element whose class is nonzero but nilpotent modulo $(f)$.
\begin{solution}
Write
\[
f=u\prod_i p_i^{e_i}
\]
with distinct irreducibles.  If some $e_i\ge2$, then $g=f/p_i$ is not in $(f)$ but a
sufficient power of $g$ is divisible by $f$, so $\bar g$ is a nonzero nilpotent in $R/(f)$.
Conversely, if all $e_i=1$ and $g^N\in(f)$, unique factorization forces every $p_i$ to
divide $g$, hence $f\mid g$ and $g\in(f)$.  Thus $(f)$ is radical.
The polynomial $xy$ is squarefree, so its quotient is reduced.  The cusp polynomial
$x^2-y^3$ is irreducible (hence squarefree) in the usual polynomial ring and gives a
reduced quotient.  The polynomial $x^2y$ has a repeated factor $x$ and gives a nonreduced
quotient.
\end{solution}
\paragraph{What happened mathematically.}
Repeated components in the equation produce nilpotent multiplicity even when the support is
unchanged.
\paragraph{Extensions.}
Relate the one-variable case to $\gcd(f,f')$ when the derivative criterion is applicable.
\paragraph{Provenance.}
New canonical problem connecting radical ideals from Volume V with affine hypersurface geometry.
